\documentclass[11pt,a4paper]{article}
\usepackage[margin=1in]{geometry}
\usepackage{amsmath,amssymb,bm}
\usepackage{hyperref}
\usepackage{graphicx}
\usepackage{booktabs}
\usepackage{enumitem}

\newcommand{\vect}[1]{\bm{#1}}
\newcommand{\unit}[1]{\,\mathrm{#1}}
\newcommand{\fm}{\unit{fm}}
\newcommand{\GeV}{\unit{GeV}}
\newcommand{\MeV}{\unit{MeV}}

\title{\texttt{genQE\_FSI} --- Two-Nucleon GCF Generator\\
       with GENIE Intranuclear Cascade\\[4pt]
       \large Physics \& Pipeline Documentation}
\author{GCF with 3N Project}
\date{\today}

\begin{document}
\maketitle
\tableofcontents
\newpage

%======================================================================
\section{Overview}
%======================================================================

\texttt{genQE\_FSI} wraps the two-nucleon GCF quasi-elastic generator
(\texttt{genQE}) with a \textbf{Final State Interaction (FSI)} layer.
After the GCF vertex produces a correlated nucleon pair, both the
struck (``lead'') and partner (``recoil'') nucleons are independently
transported through the residual nucleus using GENIE's
\textbf{Intranuke2018} intranuclear cascade (INC).

The reaction is
\begin{equation}
  e + A \;\longrightarrow\; e' + N_{\text{lead}}' + N_{\text{rec}}'
                           + X + (A{-}2)^{*},
\end{equation}
where primed nucleons denote post-FSI momenta and $X$ represents any
additional secondaries (pions, ejected nucleons) produced during
re-scattering.

Two cascade models are available:
\begin{itemize}[nosep]
  \item \textbf{hN} (\texttt{HNIntranuke2018}): full hadron--nucleon
        cascade with explicit angular distributions.  More accurate.
  \item \textbf{hA} (\texttt{HAIntranuke2018}): effective hadron--nucleus
        cascade using data-driven total cross sections.  Faster
        (${\sim}\,5$--$10\times$), suitable for systematics studies.
\end{itemize}

%======================================================================
\section{Pre-FSI Stage: GCF Event Generation}
%======================================================================

The GCF vertex is identical to \texttt{genQE} (see its \texttt{PHYSICS.tex}).
It produces:
\begin{itemize}[nosep]
  \item Lead and recoil nucleon 4-momenta
        $(\vect{p}_{\text{lead}},\,\vect{p}_{\text{rec}})$
  \item Pair types (PDG codes)
  \item Scattered electron 4-momentum $\vect{k}$
  \item Event weight $w$ from the GCF spectral function and $e$--$N$
        cross section
\end{itemize}

At this stage the nucleons have momenta but \emph{no spatial position} ---
the GCF is a momentum-space formalism.

%======================================================================
\section{SRC Position Sampling}
\label{sec:position}
%======================================================================

FSI requires a starting position inside the nucleus.  Since the SRC pair
was formed by two nucleons overlapping at short range, the probability of
finding the pair at radius $r$ scales as the \emph{square} of the
single-nucleon density:
\begin{equation}\label{eq:pos}
  P(r)\,dr \;\propto\; \rho^{2}(r)\;4\pi r^{2}\,dr.
\end{equation}

\subsection{Density models}

GENIE's \texttt{genie::utils::nuclear::Density($r$,\,$A$)} returns the
single-nucleon density using:

\paragraph{Woods--Saxon (A $>$ 20).}
\begin{equation}
  \rho_{\text{WS}}(r) = \frac{\rho_{0}}{1 + \exp\!\bigl[(r - R)/a\bigr]},
  \qquad
  R = r_{0}\,A^{1/3},\quad a \approx 0.545\fm,
\end{equation}
with $r_{0} \approx 1.07\fm$ and $\rho_{0}$ set by normalization to
$\rho_{\text{normal}} = 0.17\fm^{-3}$.

\paragraph{Modified Gaussian ($4 < A \le 20$).}
\begin{equation}
  \rho_{\text{G}}(r) = \rho_{0}\,
    \exp\!\left[-\left(\frac{r}{R}\right)^{2}\right],
  \qquad R = r_{0}\,A^{1/3}.
\end{equation}

Parameters are fitted to electron-scattering charge-density data
(de Vries \textit{et al.}, At.\ Data Nucl.\ Data Tables \textbf{36}, 1987).

\subsection{Rejection sampling algorithm}

The radial coordinate is sampled via rejection:
\begin{enumerate}
  \item Compute $R_{\max} = 2.5 \times 1.4 \times A^{1/3}\fm$.
  \item Find the envelope $f_{\max} = 1.01\,\max_{r}\bigl[\rho^{2}(r)\,r^{2}\bigr]$
        by scanning 200 grid points (cached per $A$).
  \item Draw $r \sim \text{Uniform}(0,\,R_{\max})$.
  \item Accept if $U \cdot f_{\max} \le \rho^{2}(r)\,r^{2}$
        with $U \sim \text{Uniform}(0,1)$.
\end{enumerate}

An isotropic direction $(\cos\theta,\,\phi)$ is drawn uniformly.
Both SRC nucleons are placed at the \emph{same} spatial position
$\vect{x}_{\text{SRC}}$, reflecting their short-range overlap.

\subsection{Caching}

The envelope $f_{\max}$ and $R_{\max}$ are stored as statics and
recomputed only when $A$ changes, avoiding ${\sim}\,200$ density
evaluations per event.

%======================================================================
\section{GENIE Intranuclear Cascade --- General Framework}
\label{sec:cascade}
%======================================================================

The remainder of this document describes how GENIE's hN cascade works in
detail.  The hA model follows similar logic with data-driven effective
cross sections.

\subsection{Event record construction}
\label{sec:evrec}

GENIE's cascade operates on a \texttt{GHepRecord} (event record).  We
construct a minimal record to transport a single nucleon:

\begin{center}
\begin{tabular}{clll}
\toprule
Slot & PDG & Status & Purpose \\
\midrule
0 & $e^{-}$ (11) & \texttt{kIStInitialState} &
    Fake probe; triggers \texttt{kGMdLeptonNucleus} mode \\
1 & Nucleus ($A$,$Z$) & \texttt{kIStInitialState} &
    Target; used by \texttt{SetTrackingRadius} \\
2 & Nucleus ($A$,$Z$) & \texttt{kIStStableFinalState} &
    Remnant nucleus \\
3 & Nucleon & \texttt{kIStHadronInTheNucleus} &
    Particle to transport; carries sampled position \\
\bottomrule
\end{tabular}
\end{center}

\paragraph{Why a fake electron?}
GENIE determines the ``event generation mode'' from the PDG of slot~0.
If it sees a lepton ($e$, $\mu$, $\nu$), it enters
\texttt{kGMdLeptonNucleus} mode, which \emph{skips}
\texttt{GenerateVertex()} --- the routine that would otherwise overwrite
our sampled position with a random point on the nuclear surface.

The fake electron is set to $E = 100\GeV$ to avoid triggering energy
sanity checks in the hA model.  It is never transported (only particles
with status \texttt{kIStHadronInTheNucleus} are stepped).

\subsection{Tracking radius}

Before the cascade begins, GENIE sets a tracking sphere:
\begin{equation}
  R_{\text{track}} = f_{NR} \times R_{0} \times A^{1/3},
  \qquad
  R_{0} = 1.4\fm,\quad f_{NR} = 3.
\end{equation}
A hadron is considered ``inside the nucleus'' while
$|\vect{x}| < R_{\text{track}} + \delta r$, where
$\delta r = 0.05\fm$ (one step).

\subsection{Stepping}

The cascade is semi-classical: the nucleon is stepped through the
nucleus in fixed increments of
\begin{equation}
  \delta r = 0.05\fm.
\end{equation}

At each step the particle's position is advanced along its momentum
direction.  A boundary check prevents it from flying more than $1\fm$
past the nuclear surface.

%======================================================================
\section{Mean Free Path and Interaction Probability}
%======================================================================

At each step the interaction probability is determined by the
\textbf{mean free path}
\begin{equation}\label{eq:mfp}
  \lambda(r) = \frac{1}{\rho(r)\,\sigma_{\text{tot}}},
\end{equation}
where $\rho(r)$ is the local nuclear density and $\sigma_{\text{tot}}$ is
the total hadron--nucleon cross section (in $\fm^{2}$).

\subsection{Nuclear density with de Broglie correction}

The density entering Eq.~\eqref{eq:mfp} is the full nuclear density
$A\,\rho(r)$ (i.e., number density, not single-nucleon density), with an
empirical \emph{de Broglie ring correction} that accounts for the finite
size of the interaction region:
\begin{equation}
  \rho_{\text{eff}}(r) = A\;\rho\!\left(r,\;A,\;\text{ring}\right),
  \qquad
  \text{ring} = n \times \frac{\hbar c}{p},
\end{equation}
where $p$ is the hadron's momentum and $n$ depends on $A$ and particle
type:
\begin{itemize}[nosep]
  \item Pions: $n = \delta R_{\pi}$ (for $A > 20$, doubled for light nuclei)
  \item Nucleons: $n = \delta R_{N}$ (similar)
\end{itemize}
Default values: $\delta R_{\pi} \approx \delta R_{N} \approx 1.5$.

\subsection{Total cross section}

For a \textbf{proton} at kinetic energy $T$:
\begin{equation}
  \sigma_{\text{tot}}^{p}
  = \sigma_{pp}(T)\,\frac{Z}{A}
  + \sigma_{pn}(T)\,\frac{A-Z}{A}
  + \sigma_{\text{Coulomb}},
\end{equation}
where $\sigma_{pp}$, $\sigma_{pn}$ are parameterized from data
(tabulated splines in \texttt{INukeHadroData2018}, clamped to
$T \in [10,\,1000]\MeV$).  A Coulomb correction (Cohen formula with
penetrability factor) is added for protons.

An empirical low-energy cutoff sets $\sigma_{\text{tot}} = 0$ if
$T < A^{0.2} \times 12\MeV$.

For \textbf{pions} at $T \ge 350\MeV$:
\begin{equation}
  \sigma_{\text{tot}}^{\pi}
  = \sigma_{\pi p}(T)\,\frac{Z}{A}
  + \sigma_{\pi n}(T)\,\frac{A-Z}{A}.
\end{equation}
For $T < 350\MeV$ the Oset model is used instead
(Section~\ref{sec:oset}).

\subsection{Pandharipande--Pieper medium corrections}
\label{sec:pp}

For nucleon--nucleon scattering, the free-space cross section is modified
by an in-medium correction factor:
\begin{equation}
  \sigma_{\text{tot}}^{NN,\text{med}}
  = \sigma_{\text{tot}}^{NN,\text{free}}
    \times f_{\text{PP}}(\rho,\,T),
\end{equation}
where $f_{\text{PP}} \in [0.3,\,0.8]$ typically, accounting for:

\paragraph{Pauli blocking.}
The fraction of NN scattering final states that are \emph{above} the local
Fermi surface.  The local Fermi momentum is
\begin{equation}
  k_{F,p} = \left(\frac{3\pi^{2}\,\rho\,Z}{A}\right)^{\!1/3}\!\!\frac{1}{\hbar c},
  \qquad
  k_{F,n} = \left(\frac{3\pi^{2}\,\rho\,(A-Z)}{A}\right)^{\!1/3}\!\!\frac{1}{\hbar c}.
\end{equation}

\paragraph{In-medium effective mass.}
The nucleon acquires a momentum-dependent effective mass
\begin{equation}
  m^{*}(\rho,\,k^{2})
  = m_{N}\,\frac{(\Lambda^{2} + k^{2})^{2}}
               {(\Lambda^{2} + k^{2})^{2} - 2\Lambda^{2}\,B\,m_{N}},
\end{equation}
with $\Lambda(\rho) = \lambda_{0} + \lambda_{1}\,\rho$ and
$B(\rho) = \beta_{1}\,\rho$ parameterized from nuclear-matter
calculations (Pandharipande \& Pieper, Phys.\ Rev.\ C \textbf{45}, 791, 1992).

The correction is pre-tabulated on a $200 \times 17$ grid in
$(T,\,\rho)$ and bilinearly interpolated.

\subsection{Interaction sampling}

At each step, an interaction distance is drawn from an exponential:
\begin{equation}
  d = -\lambda\,\ln U, \qquad U \sim \text{Uniform}(0,1).
\end{equation}
If $d < \delta r = 0.05\fm$, the hadron interacts; otherwise it
continues to the next step.

%======================================================================
\section{Hadron Fates}
\label{sec:fates}
%======================================================================

When an interaction occurs, a \textbf{fate} is selected by drawing from
the relative cross sections.  Six outcomes are possible:

\subsection{Elastic scattering (\texttt{kIHNFtElas})}

The hadron scatters elastically off a single bound nucleon.  No particle
species change; only momenta are redistributed.

\paragraph{Kinematics.}
The scattering angle $\cos\theta_{\text{CM}}$ is sampled from
parameterized NN angular distributions (tabulated in
\texttt{INukeHadroData2018}).  Full relativistic two-body kinematics are
applied:
\begin{enumerate}
  \item Boost to the CM frame of the hadron--nucleon system.
  \item Compute outgoing CM momenta:
        \begin{equation}
          E_{3}^{\text{CM}} = \frac{s + m_{3}^{2} - m_{4}^{2}}{2\sqrt{s}},
          \qquad
          |\vect{p}_{3}^{\text{CM}}| = \sqrt{(E_{3}^{\text{CM}})^{2} - m_{3}^{2}}.
        \end{equation}
  \item Rotate by sampled $(\theta_{\text{CM}},\,\phi_{\text{CM}})$ with
        $\phi$ uniform in $[0,2\pi)$.
  \item Boost back to the lab frame.
\end{enumerate}

The struck nucleon's Fermi momentum is sampled from the local Fermi gas
(when \texttt{fDoFermi} = true, default).  The remnant nucleus 4-momentum
is updated accordingly.

\subsection{Charge exchange (\texttt{kIHNFtCEx})}

The hadron exchanges charge with a bound nucleon:
\begin{align}
  \pi^{+} + n &\to \pi^{0} + p, &
  \pi^{-} + p &\to \pi^{0} + n, \\
  \pi^{0} + p &\to \pi^{+} + n, &
  \pi^{0} + n &\to \pi^{-} + p.
\end{align}
For nucleons, charge exchange corresponds to $p + n \to n + p$ (with
different momentum sharing).  The kinematics are identical to elastic
scattering but with different outgoing particle masses and PDG codes.

\paragraph{Availability.}
Charge exchange requires the appropriate partner nucleon to be present
in the remnant (e.g., $\pi^{+}$~CEx needs at least one neutron).  If
unavailable, this fate's weight is set to zero.

\subsection{Inelastic / pion production (\texttt{kIHNFtInelas})}

At sufficiently high energy, the collision produces one or more pions:
\begin{equation}
  N + N \to N + N + n\pi, \qquad
  \pi + N \to N + m\pi.
\end{equation}
The final state is generated using GENIE's hadronization routines, which
sample multiplicity and kinematics from parameterized inclusive data.
Each produced pion is itself subject to further cascade steps.

\subsection{Absorption (\texttt{kIHNFtAbs})}

The hadron is absorbed by a correlated nucleon pair (typically a
deuteron-like $pn$ system):
\begin{align}
  \pi^{+} + (pn) &\to p + p, \\
  \pi^{-} + (pn) &\to n + n, \\
  \pi^{0} + (pn) &\to p + n.
\end{align}
This requires \emph{both} a proton and a neutron in the remnant.

\paragraph{Kinematics.}
Two target nucleons are sampled with Fermi momenta (scaled by an
empirical factor $f_{\text{Fermi}} = 0.6$).  Two-body kinematics are
applied in the $\pi + (NN)$ CM frame.  The outgoing nucleons are
Pauli-blocked: if either falls below $k_{F}$, the absorption is rejected
and re-sampled.

The remnant loses two nucleons: $A_{\text{rem}} \to A_{\text{rem}} - 2$,
$Z_{\text{rem}} \to Z_{\text{rem}} - \Delta Z$.

\subsection{Compound nucleus (\texttt{kIHNFtCmp})}

At very low kinetic energies ($T \lesssim 30\MeV$ for $A > 5$), a
nucleon may be captured into the remnant nucleus, forming a compound
system.  The compound nucleus subsequently decays via pre-equilibrium
emission (a fast nucleon is ejected carrying a fraction of the
excitation energy) or statistical equilibrium decay.

\subsection{No interaction (\texttt{kIHNFtNoInteraction})}

The hadron escapes the nucleus without any interaction.  This is the
outcome when $d > \delta r$ at every step until
$|\vect{x}| > R_{\text{track}}$.

\subsection{Fate selection}

Each fate's probability is proportional to its partial cross section.
For nucleons:
\begin{equation}
  P(\text{fate})
  = \frac{\sigma_{\text{fate}}(T,\,A,\,Z)}
         {\sigma_{\text{elas}} + \sigma_{\text{inel}} + \sigma_{\text{cmp}}}.
\end{equation}
For pions (high energy):
\begin{equation}
  P(\text{fate})
  = \frac{\sigma_{\text{fate}}}
         {\sigma_{\text{CEx}} + \sigma_{\text{elas}} + \sigma_{\text{inel}}
          + \sigma_{\text{abs}}}.
\end{equation}

Each partial cross section can be scaled by a tunable
``fate weight'' factor (e.g., \texttt{fNucAbsFac} for absorption).
For $\pi^{0}$ absorption, an additional isospin factor of 0.665 is
applied.

A cumulative distribution is formed and a uniform random number selects
the fate.

%======================================================================
\section{Oset Model for Low-Energy Pions}
\label{sec:oset}
%======================================================================

For pion kinetic energies $T_{\pi} < 350\MeV$, the cascade uses the
\textbf{Oset model} (Salcedo \textit{et al.}, Nucl.\ Phys.\ A
\textbf{484}, 557, 1988), which provides density-dependent pion--nucleon
cross sections dominated by the $\Delta(1232)$ resonance.

\subsection{$\Delta$ propagator}

The $\pi N$ invariant mass is
\begin{equation}
  M_{\text{inv}} = \sqrt{m_{\pi}^{2} + 2m_{N}(T_{\pi} + m_{N}) + m_{N}^{2}}.
\end{equation}
The $\Delta$ propagator denominator is
\begin{equation}
  |D_{\Delta}|^{2}
  = (M_{\text{inv}} - m_{\Delta})^{2} + (\Gamma_{\text{red}} + \Sigma)^{2},
\end{equation}
where $\Gamma_{\text{red}}$ is the reduced (Pauli-blocked) $\Delta$ decay
width and $\Sigma$ is the in-medium self-energy.

\subsection{Reduced width}

The free-space $\Delta \to \pi N$ width at invariant mass $M_{\text{inv}}$ is
\begin{equation}
  \Gamma_{\text{free}}
  = \frac{g^{2}}{12\pi}\,f_{\text{coupl}}\,\frac{m_{N}\,p_{\text{CM}}^{3}}{M_{\text{inv}}},
  \qquad
  p_{\text{CM}} = \frac{p_{\pi}\,m_{N}}{M_{\text{inv}}}.
\end{equation}
This is multiplied by a \textbf{Pauli reduction factor}:
\begin{equation}
  R_{\text{Pauli}}
  = \begin{cases}
      0 & \mu_{0} < -1 \\[4pt]
      \dfrac{\mu_{0}^{3} + \mu_{0} + 2}{4} & -1 \le \mu_{0} \le 1 \\[4pt]
      1 & \mu_{0} > 1
    \end{cases}
\end{equation}
with
\begin{equation}
  \mu_{0}
  = \frac{E_{\Delta}\,E_{N}^{\text{CM}} - E_{F}\,M_{\text{inv}}}
         {p_{\pi}\,p_{\text{CM}}},
  \qquad
  E_{F} = \sqrt{k_{F}^{2} + m_{N}^{2}},
  \quad
  k_{F} = \left(\frac{3\pi^{2}\rho}{2}\right)^{\!1/3}\!\!\hbar c.
\end{equation}

\subsection{Self-energy}

The in-medium $\Delta$ self-energy $\Sigma$ contains absorption terms
($NN$ and $3N$) parameterized by Oset coefficients that depend on
density and $\pi N$ kinematics.

\subsection{Cross sections}

The total, quasi-elastic, charge-exchange, and absorption cross sections
are decomposed into $s$-wave and $p$-wave contributions, with
isospin-dependent coefficients for each pion charge state.

The fate fractions are:
\begin{align}
  f_{\text{abs}} &= \sigma_{\text{abs}} / \sigma_{\text{tot}}, \\
  f_{\text{CEx}} &= \sigma_{\text{CEx}} / \sigma_{\text{tot}}, \\
  f_{\text{elas}} &= 1 - f_{\text{abs}} - f_{\text{CEx}}.
\end{align}

%======================================================================
\section{Post-FSI Processing}
%======================================================================

\subsection{Nucleon selection}

After GENIE's \texttt{ProcessEventRecord()} returns, the stable
descendants of the transported nucleon (slot~3) are examined.  The
\textbf{highest-energy nucleon} among them is selected as the outgoing
lead (or recoil) nucleon.  Its PDG code may have changed (charge
exchange), and its momentum is modified by the cascade.

\subsection{Secondary collection}

All other stable descendants are stored as \texttt{FSISecondary} objects,
tagged with:
\begin{itemize}[nosep]
  \item \texttt{parentRole}: 0 (from lead transport) or 1 (from recoil)
  \item \texttt{pdg}: particle type
  \item \texttt{p4}: 4-momentum
  \item \texttt{rescatterCode}: GENIE's internal fate code (for reweighting)
\end{itemize}

\subsection{Residual nucleus}

The residual nucleus used for transport is computed by removing the SRC
pair from the full nucleus:
\begin{equation}
  Z_{\text{res}} = Z - n_{p}, \qquad
  A_{\text{res}} = A - 2,
\end{equation}
where $n_{p}$ is the number of protons in the pair.  For ultra-light
residuals ($A_{\text{res}} < 3$), the full nucleus is used instead to
avoid pathological density profiles.

\subsection{Pauli blocking (post-FSI)}

After both nucleons have been transported, a final Pauli-blocking check
is applied:
\begin{equation}
  \text{if}\quad |\vect{p}_{\text{lead}}'| < k_{F}
  \quad\text{or}\quad |\vect{p}_{\text{rec}}'| < k_{F},
  \quad\Longrightarrow\quad w = 0.
\end{equation}
Default Fermi momentum: $k_{F} = 0.220\GeV/c$ (tunable via
\texttt{SetFSITuning}).

\subsection{Event classification}

Each event is classified by:
\begin{itemize}[nosep]
  \item \textbf{Charge exchange}: lead or recoil changed PDG
        ($p \leftrightarrow n$)
  \item \textbf{Elastic-like}: momentum changed, no charge exchange,
        no non-nucleon secondaries produced
  \item \textbf{Pion counts}: $n_{\pi^{+}}$, $n_{\pi^{-}}$, $n_{\pi^{0}}$
        tallied separately for lead and recoil transports
\end{itemize}

%======================================================================
\section{Full Pipeline Summary}
%======================================================================

For each Monte Carlo trial:
\begin{enumerate}
  \item \textbf{GCF vertex}: Select pair type, sample CM and relative
        momenta, compute electron vertex, evaluate $\sigma_{eN}$
        (identical to \texttt{genQE}).
  \item \textbf{Position sampling}: Draw $\vect{x}_{\text{SRC}}$ from
        $P(r) \propto \rho^{2}(r)\,r^{2}$ via rejection sampling.
  \item \textbf{Compute residual nucleus}: $A_{\text{res}} = A - 2$,
        $Z_{\text{res}} = Z - n_{p}$.
  \item \textbf{Transport lead nucleon}: Build GENIE event record with
        nucleon at $\vect{x}_{\text{SRC}}$; run
        \texttt{ProcessEventRecord()}; extract highest-energy nucleon and
        secondaries.
  \item \textbf{Transport recoil nucleon}: Same procedure, same starting
        position, same residual nucleus.
  \item \textbf{Classify}: Detect charge exchange, elastic-like events,
        count pions.
  \item \textbf{Pauli blocking}: Reject event if either final nucleon
        momentum $< k_{F}$.
  \item \textbf{Store}: Write event to ROOT tree if $w > 0$.
\end{enumerate}

%======================================================================
\section{Configuration and Tuning Parameters}
%======================================================================

\begin{center}
\begin{tabular}{llll}
\toprule
Parameter & Default & Setter & Description \\
\midrule
FSI model & hN & \texttt{SetFSIModel(kHN2018)} & hN or hA cascade \\
FSI on/off & on & \texttt{EnableFSI(true)} & Bypass cascade \\
$k_{F}$ & 220 MeV/$c$ & \texttt{SetFSITuning(220)} & Pauli blocking threshold \\
$\delta r$ & 0.05 fm & (GENIE internal) & Cascade step size \\
$R_{0}$ & 1.4 fm & (GENIE internal) & Nuclear radius coefficient \\
$f_{NR}$ & 3.0 & (GENIE internal) & Tracking radius multiplier \\
$\delta R_{\pi}$ & ${\sim}1.5$ & (GENIE config) & Pion de Broglie ring \\
$\delta R_{N}$ & ${\sim}1.5$ & (GENIE config) & Nucleon de Broglie ring \\
NN medium corr. & on & \texttt{INUKE-XsecNNCorr} & Pandharipande--Pieper \\
Oset model & on & \texttt{HNINUKE-UseOset} & Low-energy pion cross sections \\
\bottomrule
\end{tabular}
\end{center}

GENIE configuration is loaded from XML files in\\
\texttt{config/genie\_override/} (e.g., \texttt{HNIntranuke2018.xml},
\texttt{Messenger.xml}).

%======================================================================
\section{References}
%======================================================================

\begin{enumerate}[label={[\arabic*]}]
  \item GENIE Collaboration, ``Physics and User Manual,''
        \url{https://genie-mc.org}, GENIE v3.06.02.
  \item C.~Andreopoulos \textit{et al.},
        ``The GENIE Neutrino Monte Carlo Generator,''
        Nucl.\ Instrum.\ Meth.\ A \textbf{614}, 87 (2010).
  \item L.~L.~Salcedo \textit{et al.},
        ``Computer simulation of inclusive pion nuclear reactions,''
        Nucl.\ Phys.\ A \textbf{484}, 557 (1988).
  \item V.~R.~Pandharipande and S.~C.~Pieper,
        ``Nuclear transparency to intermediate-energy nucleons from
        $(e,e'p)$ reactions,''
        Phys.\ Rev.\ C \textbf{45}, 791 (1992).
  \item H.~de~Vries, C.~W.~de~Jager, and C.~de~Vries,
        ``Nuclear charge-density-distribution parameters from elastic
        electron scattering,''
        At.\ Data Nucl.\ Data Tables \textbf{36}, 495 (1987).
  \item R.~Weiss \textit{et al.},
        ``Generalized nuclear contacts and momentum distributions,''
        Phys.\ Rev.\ C \textbf{92}, 054311 (2015).
  \item R.~Cruz-Torres \textit{et al.},
        ``Short-range correlations and the nuclear contact formalism,''
        Phys.\ Lett.\ B \textbf{785}, 304 (2018).
\end{enumerate}

\end{document}
